\chapter{Trabalhos Relacionados}
\label{cap:trabalhos-relacionados}

O mercado de aplicações voltadas à gestão pecuária possui diversas soluções destinadas ao controle de rebanho, produção, estoque e acompanhamento sanitário dos animais. Entretanto, apesar da ampla quantidade de funcionalidades disponíveis, muitos desses sistemas apresentam interfaces complexas, excesso de informações simultâneas e fluxos operacionais extensos, o que pode dificultar sua utilização por pequenos produtores rurais ou usuários com menor familiaridade tecnológica.

Neste capítulo são analisados aplicativos com proposta semelhante à do sistema Curral, destacando suas principais funcionalidades, características de interface e limitações observadas durante o uso. A análise busca identificar lacunas existentes nas soluções atuais e justificar a proposta de uma aplicação com foco em simplicidade operacional, acessibilidade e rapidez no registro de informações.

\section{JetBov de Campo}
\label{sec:jetbov}

O aplicativo JetBov de Campo é uma solução voltada ao gerenciamento de propriedades pecuárias, oferecendo funcionalidades relacionadas ao controle de rebanho, manejo sanitário, reprodução e gestão financeira. A aplicação possui um conjunto abrangente de recursos e funcionalidades destinadas à administração de propriedades rurais de médio e grande porte.

Durante a utilização do aplicativo, observou-se que parte das funcionalidades exige múltiplas etapas para conclusão de operações simples, como registro de manejo ou atualização de informações dos animais. Além disso, a interface apresenta grande quantidade de informações simultaneamente, o que pode aumentar a curva de aprendizado para usuários iniciantes.

	\begin{figure}[h!]
		\centering
		\Caption{\label{fig:jetbov-home} Tela inicial do aplicativo JetBov de Campo}	
		\UNIFORfig{}{
			\fbox{\includegraphics[width=8cm]{figuras/jetbov-home}}
		}{
			\Fonte{Elaborado pelo autor}
		}	
	\end{figure}

Outro aspecto identificado foi a elevada quantidade de opções presentes nos menus de navegação. Embora isso amplie a capacidade de gerenciamento do sistema, a organização visual pode dificultar a localização rápida de funcionalidades específicas durante o uso cotidiano.

\section{BovPilot}
\label{sec:bovpilot}

O BovPilot é um aplicativo direcionado à gestão de rebanhos e controle zootécnico, permitindo o registro de informações relacionadas à identificação dos animais, produção e eventos sanitários. A solução apresenta funcionalidades relevantes para acompanhamento produtivo e histórico do rebanho.

Durante a análise da aplicação, verificou-se que determinadas operações exigem navegação entre múltiplas telas e formulários. Esse comportamento pode tornar o processo de inserção de dados mais lento em situações de uso frequente no ambiente rural.

	\begin{figure}[h!]
		\centering
		\Caption{\label{fig:bovpilot-cadastro} Fluxo de cadastro de informações no aplicativo BovPilot}	
		\UNIFORfig{}{
			\fbox{\includegraphics[width=8cm]{figuras/bovpilot-cadastro}}
		}{
			\Fonte{Elaborado pelo autor}			
		}	
	\end{figure}

Também foi identificado o uso recorrente de terminologias técnicas e elementos visuais densos, o que pode dificultar a adaptação inicial de usuários sem experiência prévia com sistemas de gestão pecuária digital.

\section{Gado Certo}
\label{sec:gado-certo}

O aplicativo Gado Certo possui funcionalidades relacionadas ao gerenciamento de animais, controle financeiro e acompanhamento produtivo. A solução busca centralizar informações da propriedade rural em um único ambiente digital.

Durante a utilização do aplicativo, observou-se que algumas telas apresentam excesso de campos e informações simultaneamente. Em dispositivos móveis com telas reduzidas, esse comportamento pode comprometer a legibilidade e a velocidade de operação do sistema.

	\begin{figure}[h!]
		\centering
		\Caption{\label{fig:gado-certo-interface} Interface de gerenciamento do aplicativo Gado Certo}	
		\UNIFORfig{}{
			\fbox{\includegraphics[width=8cm]{figuras/gado-certo-interface}}
		}{
			\Fonte{Elaborado pelo autor}
		}	
	\end{figure}

Além disso, determinadas funcionalidades exigem diversos passos até a conclusão da operação desejada, aumentando o tempo necessário para registros simples realizados diariamente pelo produtor rural.

\section{Gerente de Pecuária}
\label{sec:gerente-pecuaria}

O aplicativo Gerente de Pecuária oferece funcionalidades voltadas ao acompanhamento do rebanho, controle sanitário, estoque e gerenciamento produtivo. A solução possui abordagem semelhante à de outros aplicativos do segmento, concentrando diferentes módulos administrativos em uma única plataforma.

Durante a análise do sistema, verificou-se que a organização visual das funcionalidades pode dificultar a identificação rápida das operações mais utilizadas no cotidiano da propriedade rural. Além disso, a navegação apresenta quantidade significativa de opções e telas intermediárias.

	\begin{figure}[h!]
		\centering
		\Caption{\label{fig:gerente-pecuaria} Tela de gerenciamento do aplicativo Gerente de Pecuária}	
		\UNIFORfig{}{
			\fbox{\includegraphics[width=8cm]{figuras/gerente-pecuaria}}
		}{
			\Fonte{Elaborado pelo autor}
		}	
	\end{figure}

Observou-se ainda que parte das interfaces prioriza densidade de funcionalidades em detrimento da simplicidade operacional, podendo gerar dificuldades para usuários que necessitam de registros rápidos em ambientes externos ou durante atividades de campo.

\section{Análise Comparativa}
\label{sec:analise-comparativa}

Os aplicativos analisados apresentam funcionalidades relevantes para gestão pecuária e demonstram evolução significativa na digitalização de processos rurais. Entretanto, foi possível identificar características recorrentes entre as soluções avaliadas, como interfaces densas, elevada quantidade de funcionalidades simultâneas e fluxos operacionais extensos.

	\begin{quadro}[h!]	
		\centering
		\Caption{\label{qua:comparativo-apps} Comparação entre os aplicativos analisados}		
		\UNIFORqua{}{
			\begin{tabular}{|p{3cm}|p{4cm}|p{5cm}|}
				\hline
				Aplicativo & Principais funcionalidades & Limitações observadas \\
				\hline
				JetBov de Campo & Gestão completa do rebanho e manejo sanitário & Interface com grande quantidade de informações e menus extensos \\
				\hline
				BovPilot & Controle zootécnico e histórico animal & Fluxos de cadastro com múltiplas etapas \\
				\hline
				Gado Certo & Gestão financeira e produtiva & Excesso de campos em telas móveis \\
				\hline
				Gerente de Pecuária & Controle integrado da propriedade & Navegação extensa e elevada densidade visual \\
				\hline
			\end{tabular}
		}{
			\Fonte{Elaborado pelo autor}
		}
	\end{quadro}

A partir das análises realizadas, observa-se a existência de uma lacuna relacionada à simplicidade de uso e à acessibilidade operacional em aplicações voltadas ao pequeno produtor rural. Nesse contexto, o aplicativo Curral propõe uma abordagem centrada em fluxos simplificados, navegação reduzida e rapidez no registro das informações, buscando diminuir barreiras de utilização e facilitar a adoção tecnológica no ambiente rural.

	\begin{quadro}[h!]	
		\centering
		\Caption{\label{qua:diferenciais-curral} Diretrizes adotadas pelo aplicativo Curral}		
		\UNIFORqua{}{
			\begin{tabular}{|p{5cm}|p{8cm}|}
				\hline
				Diretriz & Objetivo \\
				\hline
				Interface simplificada & Reduzir a complexidade visual durante o uso \\
				\hline
				Fluxos curtos de cadastro & Diminuir o tempo necessário para inserção de informações \\
				\hline
				Navegação objetiva & Facilitar a localização das funcionalidades principais \\
				\hline
				Operação mobile-first & Melhorar a utilização em ambientes externos e rurais \\
				\hline
				Entrada rápida de dados & Permitir registros cotidianos com menor quantidade de etapas \\
				\hline
			\end{tabular}
		}{
			\Fonte{Elaborado pelo autor}
		}
	\end{quadro}

Dessa forma, os trabalhos relacionados analisados contribuem para identificar limitações presentes nas soluções atuais e fundamentam as decisões arquiteturais e de interface adotadas no desenvolvimento do aplicativo Curral.
\section{Análise de Soluções Existentes e Lacunas de Mercado}
\label{sec:lacunas-mercado}

A análise das soluções atualmente disponíveis para gestão pecuária no cenário brasileiro evidencia um contraste entre a existência de ferramentas digitais especializadas e sua aplicabilidade prática no contexto da agricultura familiar. Embora exista uma variedade de sistemas voltados ao gerenciamento rural, parte dessas soluções apresenta limitações relacionadas à acessibilidade operacional, adaptação ao pequeno produtor e simplicidade de utilização.

Durante a pesquisa por ferramentas de gestão pecuária, observou-se a presença tanto de sistemas consolidados historicamente quanto de aplicações comerciais modernas. Entretanto, muitas dessas soluções não atendem plenamente às necessidades identificadas no contexto deste trabalho, especialmente em propriedades rurais de pequeno porte e em cenários de uso cotidiano no interior do Ceará.

Um exemplo relevante é o sistema Gerenpec, desenvolvido pela Embrapa. Apesar de sua importância histórica no planejamento de propriedades pecuárias, as principais referências relacionadas à ferramenta remontam ao início dos anos 2000, período anterior à popularização de dispositivos móveis e das práticas associadas ao conceito de Agro 4.0. Dessa forma, sua abordagem apresenta foco predominante em planejamento macro da produção, diferindo da necessidade observada neste trabalho, que envolve acompanhamento diário e individualizado dos animais.

Além disso, observou-se que muitos sistemas tradicionais priorizam funcionalidades administrativas e estratégicas relacionadas ao rebanho de forma global, enquanto pequenos produtores frequentemente necessitam de mecanismos simples e rápidos para registro de vacinação, controle sanitário, acompanhamento de peso e produção individual dos animais.

Por outro lado, aplicativos comerciais contemporâneos, como JetBov, BovPilot, Gado Certo e Gerente de Pecuária, oferecem funcionalidades mais modernas e maior integração tecnológica. Entretanto, a análise prática dessas aplicações permitiu identificar limitações recorrentes relacionadas à complexidade operacional e à adaptação ao contexto da agricultura familiar.

As principais limitações observadas podem ser agrupadas em três categorias principais:

\begin{itemize}
    \item \textbf{Foco em propriedades de maior porte}: diversas soluções são direcionadas principalmente à pecuária de corte em larga escala, incorporando funcionalidades avançadas de gestão corporativa que nem sempre são necessárias para pequenos produtores rurais.
    
    \item \textbf{Complexidade de interface e operação}: muitos aplicativos apresentam interfaces densas, grande quantidade de menus e fluxos extensos de cadastro, exigindo maior familiaridade tecnológica por parte do usuário.
    
    \item \textbf{Ênfase na gestão coletiva do rebanho}: parte das soluções prioriza indicadores globais da fazenda, enquanto pequenos produtores frequentemente demandam acompanhamento individualizado do histórico sanitário e produtivo dos animais.
\end{itemize}

Outro aspecto relevante identificado durante a análise foi a presença de interfaces com excesso de informações simultâneas, múltiplas etapas para execução de tarefas simples e elevada densidade visual em dispositivos móveis. Em ambientes rurais, nos quais os registros frequentemente precisam ser realizados rapidamente durante atividades de campo, tais características podem dificultar a adoção tecnológica e reduzir a eficiência operacional.

Nesse contexto, observa-se a existência de uma lacuna relacionada ao desenvolvimento de soluções voltadas especificamente à realidade da agricultura familiar, com foco em simplicidade operacional, clareza visual e rapidez na entrada de dados. Essa lacuna tecnológica fundamenta a proposta do aplicativo Curral, cujo objetivo consiste em oferecer uma ferramenta acessível, intuitiva e adaptada às necessidades cotidianas do pequeno produtor rural.

Dessa forma, os trabalhos relacionados analisados não apenas demonstram o avanço das tecnologias de gestão pecuária, mas também evidenciam limitações que justificam o desenvolvimento de uma solução centrada em usabilidade, acessibilidade e gestão individualizada dos animais.